\chapter{Chapter07 — Linear CSP Algorithms}
%%% generated by the blueprint pipeline from FABL.Chapter07.LinearCSPAlgorithms
\section{Overview}

Declarations of the Lean module \texttt{FABL.Chapter07.LinearCSPAlgorithms}.

\section{Declarations}

\begin{definition}[Maximum linear predicate]
\lean{FABL.MaxLinearPredicate}
\statementsource{Analysis-of-Boolean-Functions-by-Ryan-ODonnell}{
  pdf-ffbe9c7b3306-p184-b003,
  pdf-ffbe9c7b3306-p185-b002,
  pdf-ffbe9c7b3306-p185-b003
}
\label{FABL.MaxLinearPredicate}
\leanok
\uses{FABL.𝔽₂}
For a natural number $r$, a maximum linear predicate consists of an arity $k$ with $1\le
k\le r$ and a right-hand side in the two-element field $\mathbb F_2$.
\end{definition}

\begin{definition}[Maximum linear constraint template]
\lean{FABL.maxLinearTemplate}
\statementsource{Analysis-of-Boolean-Functions-by-Ryan-ODonnell}{
  pdf-ffbe9c7b3306-p184-b003,
  pdf-ffbe9c7b3306-p185-b002,
  pdf-ffbe9c7b3306-p185-b003
}
\label{FABL.maxLinearTemplate}
\leanok
\uses{FABL.CSPTemplate, FABL.MaxLinearPredicate, FABL.𝔽₂}
For a natural number $r$, the maximum linear constraint template over $\mathbb{F}_2$ assigns
each allowed predicate an arity between $1$ and $r$, a right-hand side in the binary
field, and accepts an input tuple exactly when the sum of its entries equals that
right-hand side.
\end{definition}

\begin{definition}[Linear equation of a constraint]
\lean{FABL.CSPConstraint.toF₂LinearEquation}
\leanok
\uses{FABL.CSPConstraint, FABL.F₂LinearEquation, FABL.maxLinearTemplate}
\statementsource{Analysis-of-Boolean-Functions-by-Ryan-ODonnell}{
  pdf-ffbe9c7b3306-p185-b002,
  pdf-ffbe9c7b3306-p186-b001
}
Let $r$ be a natural number and let $V$ be a finite set with decidable equality. A
constraint over the maximum linear template of arity bound $r$ determines an equation
over the binary field whose coefficient at a variable is $1$ if that variable lies in
the constraint's scope and $0$ otherwise, and whose constant is the predicate's
right-hand side.
\end{definition}

\begin{theorem}[Evaluation of the associated linear equation]
\lean{FABL.CSPConstraint.eval_toF₂LinearEquation}
\leanok
\uses{FABL.CSPConstraint, FABL.CSPConstraint.toF₂LinearEquation, FABL.F₂LinearEquation.eval, FABL.maxLinearTemplate, FABL.𝔽₂}
\difficulty{4}
Let $r$ be a natural number, let $V$ be a finite set of variables, and let $c$ be a
constraint whose predicate is a nonempty linear equation over $\bbr/2\bbr$. For any
assignment $x:V\to\bbr/2\bbr$, the associated equation evaluates to the sum of the
assigned values of the variables in the constraint’s ordered scope.
\end{theorem}

\begin{theorem}[Constraint and equation satisfaction]
\lean{FABL.CSPConstraint.isSatisfied_toF₂LinearEquation_iff}
\leanok
\uses{FABL.CSPConstraint, FABL.CSPConstraint.Satisfied, FABL.CSPConstraint.eval_toF₂LinearEquation, FABL.CSPConstraint.toF₂LinearEquation, FABL.F₂LinearEquation.IsSatisfied, FABL.maxLinearTemplate, FABL.𝔽₂}
\difficulty{3}
Let $r$ be a natural number, let $V$ be a finite set, and let $c$ be a constraint in the
maximum-arity linear template over the two-element field $\F_2$. Thus $c$ applies a
parity equation with a nonempty scope of at most $r$ distinct variables and a right-hand
side in $\F_2$. For every assignment of elements of $\F_2$ to $V$, the associated linear
equation is satisfied if and only if $c$ is satisfied.
\end{theorem}

\begin{definition}[Linear equation rows of a CSP instance]
\lean{FABL.CSPInstance.linearRows}
\label{FABL.CSPInstance.linearRows}
\leanok
\uses{FABL.CSPConstraint.toF₂LinearEquation, FABL.CSPInstance, FABL.F₂LinearEquation, FABL.maxLinearTemplate}
Let $r$ be a natural number, let $V$ be a finite set with decidable equality, and let a
CSP instance over the maximum linear predicate template of arity at most $r$ be given.
Its linear rows are the list of linear equations obtained, in multiset-list order and
with multiplicities, by translating each constraint into its associated equation over
$\mathbb F_2$.
\end{definition}

\begin{theorem}[Linear Rows and Constraint Satisfaction]
\lean{FABL.CSPInstance.satisfiesRows_linearRows_iff}
\label{FABL.CSPInstance.satisfiesRows_linearRows_iff}
\leanok
\uses{FABL.CSPConstraint.isSatisfied_toF₂LinearEquation_iff, FABL.CSPConstraint.toF₂LinearEquation, FABL.CSPInstance, FABL.CSPInstance.FullySatisfied, FABL.CSPInstance.linearRows, FABL.F₂SatisfiesRows, FABL.maxLinearTemplate, FABL.𝔽₂}
\difficulty{4}
Let $r$ be a natural number, let $V$ be a finite variable set, and consider an instance
whose constraints are nonempty linear equations over $\bbr$'s two-element field, each
involving at most $r$ distinct variables. For any assignment of field elements to $V$,
every equation in the list of rows obtained from the constraints is satisfied if and
only if every constraint in the instance is satisfied.
\end{theorem}

\begin{definition}[Gaussian assignment]
\lean{FABL.CSPInstance.gaussianAssignment}
\statementsource{Analysis-of-Boolean-Functions-by-Ryan-ODonnell}{pdf-ffbe9c7b3306-p188-b005}
\label{FABL.CSPInstance.gaussianAssignment}
\leanok
\uses{FABL.CSPInstance, FABL.CSPInstance.linearRows, FABL.maxLinearTemplate, FABL.solveF₂Rows, FABL.𝔽₂}
Let $r$ be a natural number, let $V$ be a finite set with decidable equality, and
suppose $V$ is encodable. For an instance of the maximum linear constraint template with
variables in $V$, its Gaussian assignment is the function from $V$ to the binary field
obtained by applying the row-solving operation to the instance's linear rows.
\end{definition}

\begin{theorem}[Gaussian assignment fully satisfied of satisfiable]
\lean{FABL.CSPInstance.gaussianAssignment_fullySatisfied_of_satisfiable}
\statementsource{Analysis-of-Boolean-Functions-by-Ryan-ODonnell}{pdf-ffbe9c7b3306-p188-b005}
\label{FABL.CSPInstance.gaussianAssignment_fullySatisfied_of_satisfiable}
\leanok
\uses{FABL.CSPInstance, FABL.CSPInstance.FullySatisfied, FABL.CSPInstance.Satisfiable, FABL.CSPInstance.gaussianAssignment, FABL.CSPInstance.satisfiable_iff_exists_fullySatisfied, FABL.CSPInstance.satisfiesRows_linearRows_iff, FABL.maxLinearTemplate, FABL.solveF₂Rows_satisfies_of_exists}
\difficulty{3}
% dataset-pipeline: stub -- prose pending build_informal_statements.py
Statement pending regeneration.
\end{theorem}

\begin{theorem}[Gaussian assignment value eq one of satisfiable]
\lean{FABL.CSPInstance.gaussianAssignment_value_eq_one_of_satisfiable}
\statementsource{Analysis-of-Boolean-Functions-by-Ryan-ODonnell}{pdf-ffbe9c7b3306-p188-b005}
\label{FABL.CSPInstance.gaussianAssignment_value_eq_one_of_satisfiable}
\leanok
\uses{FABL.CSPInstance, FABL.CSPInstance.Satisfiable, FABL.CSPInstance.gaussianAssignment, FABL.CSPInstance.gaussianAssignment_fullySatisfied_of_satisfiable, FABL.CSPInstance.value, FABL.CSPInstance.value_eq_one_iff_fullySatisfied, FABL.maxLinearTemplate}
\difficulty{1}
% dataset-pipeline: stub -- prose pending build_informal_statements.py
Statement pending regeneration.
\end{theorem}

\begin{definition}[Gaussian elimination work estimate]
\lean{FABL.CSPInstance.gaussianEliminationWork}
\label{FABL.CSPInstance.gaussianEliminationWork}
\leanok
\uses{FABL.CSPInstance, FABL.CSPInstance.linearRows, FABL.cubicF₂LinearSolverBackend, FABL.maxLinearTemplate}
Let $r$ be a natural number, let $V$ be a finite set with decidable equality and an
encoding, and let a CSP instance over the maximum linear predicate template with
parameter $r$ be given. Its Gaussian elimination work is the number of constraints times
$|V|+1$, plus the work of the cubic $\mathbb F_2$ linear solver backend on the
instance's list of linear equation rows.
\end{definition}

\begin{theorem}[Bound on Gaussian elimination work]
\lean{FABL.CSPInstance.gaussianEliminationWork_le}
\label{FABL.CSPInstance.gaussianEliminationWork_le}
\leanok
\uses{FABL.CSPInstance, FABL.CSPInstance.gaussianEliminationWork, FABL.CSPInstance.linearRows, FABL.cubicF₂LinearSolverBackend, FABL.cubicF₂LinearSolverBackend_work_le, FABL.maxLinearTemplate}
\difficulty{4}
Let $r$ be a natural number, let $V$ be a finite variable set, and let an instance of
the specified maximum linear constraint system over $\bbr$'s two-element field have a
nonempty multiset of constraints in which every variable occurs. Its Gaussian
elimination work is at most
\[m(|V|+1)+|V|m(|V|+2)+|V|(|V|+1),\]
where $m$ is the number of constraints counted with multiplicity.
\end{theorem}
